\chapter{Valutazione, KPI e Analytics}
\label{g:valutazione}

\section{Valutazione (protocollo pronto, sessioni da svolgere)}
Protocollo definito (Molina et al.~\cite{molina}): think-aloud +
questionari (SUS, IMI, nucleo GEQ) su 5--8 studenti, sessioni da circa
30\,min. Sezioni \textit{metodo / profilo partecipanti / risultati /
riflessioni} da compilare dopo le sessioni.

\medskip
\noindent\fbox{\parbox{\dimexpr\linewidth-2\fboxsep-2\fboxrule\relax}{%
\textbf{DA COMPLETARE (User Evaluation da svolgere).} Nessuna sessione
\`e stata ancora svolta: nessun risultato utente viene riportato in
questo documento. Compilare metodo, profili, risultati e riflessioni solo
dopo le sessioni reali.}}

\section{KPI e Analytics}
KPI candidati: retention D1/D7, tasso di mastery al primo tentativo
(soglia 80\,\%), tempo medio per boss, carte per vittoria, SUS $\geq 70$.
Analytics locali gi\`a implementati (isolati per utente$\times$campagna,
cap 200 eventi, riepilogo ``I miei numeri'' in Settings): conteggi
\path{quiz_pass}, \path{quiz_fail}, \path{boss_win},
\path{boss_lose}, \path{reward_claim}, \path{session_start} +
\path{failCount} per topic (soglia 2 per ``Da ripassare''). Specchio Hub
(\texttt{data} nelle executions): \path{quiz_completed},
\path{claim_reward}, \path{boss_defeated}. Le soglie di successo saranno
fissate dopo la prima tornata di dati reali.
